\section{Methodology}
\label{sec:methodology}

\subsection{Data}
\label{sec:data}
We used the three collections on which S3PL was published (\cref{tab:data}),
so that both baselines could be reproduced. GBM is the massNet patient-derived
xenograft (PDX) glioblastoma dataset (PR001292)~\citep{abdelmoula2022massnet}.
CAC is the MetaboLights colorectal adenocarcinoma dataset
(MTBLS415)\checkme{acquisition paper}. RCC comes from
\citet{dill2010rcc} via Cardinal~\citep{bemis2015cardinal}. The expert masks
come from the S3PL release~\citep{weigand2026spatial}.
\begin{itemize}
  \item \textbf{GBM:} sections from the same PDX line are not independent, so
        GBM is also summarised per line.
  \item \textbf{CAC:} all sections come from one biopsy, so its comparisons
        are within one specimen.
  \item \textbf{RCC:} cancer, normal and an unannotated region (45--78\% of
        pixels) are kept as three classes, with a cancer/normal-only
        sensitivity analysis.
\end{itemize}
Four sections (GBM108\_positive, GBM22\_2, 160TopL, 40TopL) were used for
development. RCC was held out until the method was frozen.

\begin{table}[t]
\caption{Datasets.}
\label{tab:data}
\centering\small
\begin{tabular}{@{}lccc@{}}
\toprule
 & GBM & CAC & RCC \\
\midrule
Sections & 8 & 8 & 8 \\
Independent units & 4 PDX lines & 1 biopsy & 8 subjects \\
$m/z$ bins & 85{,}062 & 1{,}481 & 10{,}200 \\
Pixels per section & 2{,}071--4{,}524 & 3{,}243--4{,}761 & 972--3{,}564 \\
Classes (GMM $K$) & 2 & 3 & 3 \\
\bottomrule
\end{tabular}
\end{table}

\subsection{Pipeline}
\label{sec:pipeline}
Each section is modelled on its own spectra without labels, and its selected
peaks are scored against its own masks. The design is unsupervised and
transductive. Masks never enter fitting or ranking, but they set two
configuration values, which are disclosed below.

\textbf{Spatial-msiPL.} Following msiPL, spectra are TIC-normalised and a fully
connected VAE (one hidden layer of 512 units, 5 latent dimensions) is trained
for 100 epochs with Adam (learning rate $10^{-3}$, batch size 128), using
msiPL's categorical cross-entropy reconstruction loss plus KL
divergence~\citep{abdelmoula2021peak}.
\begin{itemize}
  \item The \emph{centre-only} arm encodes one spectrum.
  \item The \emph{uniform} arm concatenates it with the mean of its up to eight
        measured $3\times3$ neighbours.
  \item Both reconstruct the centre only.
  \item The variants are learned attention over neighbours (about 0.5\% more
        parameters)~\citep{vaswani2017attention}, $5\times5$ to $9\times9$
        windows, zeroed and shuffled neighbours, and centre masking.
\end{itemize}

\textbf{Segmentation.} A full-covariance GMM is fitted to the standardised
latent means. Its number of components $K$ equals the number of annotated
classes, a class-informed choice that BIC did not support
(\cref{sec:res-encoder}).

\textbf{Rankings.} Each ranking orders the $m/z$ bins, and the top
$K_{\text{bin}}$ are scored.
\begin{itemize}
  \item \emph{IG}~\citep{sundararajan2017axiomatic} integrates the gradient of
        the assigned GMM posterior from the section-mean spectrum to 12 sampled
        pixels per component (64 trapezoidal steps), and ranks bins by mean
        absolute attribution, balanced across components.
  \item \emph{Posterior correlation} ranks bins by the absolute Pearson
        correlation of their ion images with each posterior map, balanced the
        same way, using the same GMM.
  \item \emph{Model-free references} are mean intensity, variance, Moran's~$I$
        over Moore-neighbour edges~\citep{moran1950notes}, and random.
  \item \emph{Ceiling:} a supervised oracle that correlates with the masks.
\end{itemize}
Test~V replaces the VAE with 5-component PCA, keeping the GMM and correlation
ranking. Test~S smooths the latent means with structure-adaptive weights before
the GMM~\citep{alexandrov2011efficient}.

\textbf{Baselines.} msiPL was reproduced and tuned to its published peak-count
procedure. S3PL was run from its public repository with the released settings:
10 epochs, Adam ($10^{-2}$), batch size 16, and patch size $p$ and peaks per
patch $z$ of 3/256 (GBM), 9/256 (CAC) and 5/1024 (RCC). The released code
normalises each patch by its spatial maximum, while the paper describes TIC
normalisation, so both were tested.

\subsection{Evaluation and design}
\label{sec:evaluation}
\textbf{mSCF1.} For each class mask and threshold
$t\in\{0.3,0.4,0.5,0.6\}$, a reference set of mask-correlated bins is formed
using the S3PL evaluator's rank cutoff. The F1 score between the selected bins
and the union of the reference sets is averaged over the four
thresholds~\citep{weigand2026spatial}. mSCF1 measures agreement with the
masks, not chemical identity.

\textbf{Budgets.} All compared methods select the same number of bins,
$K_{\text{bin}}$: the tuned msiPL count for GBM and CAC, and the released S3PL
target for RCC. A peak-level scale (P1) collapses neighbouring bins of one
mean-spectrum peak, so duplicates are not credited. On CAC and RCC, P1 is
nearly singleton and acts as a second budget check.

\textbf{Predeclared rules.} Every comparison's rule was written before its
outputs, in an append-only protocol.
\begin{itemize}
  \item A difference ``helps'' if its mean is at least $+0.02$ mSCF1 and it is
        positive in at least 6/8 sections.
  \item Two methods are ``equivalent'' if the mean difference is within
        $\pm0.02$.
  \item Otherwise the label is ``mixed'', or ``incomplete'' if sections are
        missing.
\end{itemize}
All arms use training seed~1. The two core arms were retrained with seeds 2
and~3 on all sections (Test~R). Seed results are given as mean~$\pm$~SD across
seeds of the seed-level collection means.

\subsection{Statistics}
\label{sec:stats}
The labels are accompanied by paired, per-section statistics:
\begin{itemize}
  \item the exact two-sided Wilcoxon signed-rank
        test~\citep{wilcoxon1945individual};
  \item the Hodges--Lehmann (HL) median difference with an exact 95\%
        CI~\citep{hodges1963estimates};
  \item the matched-pairs rank-biserial correlation $r$.
\end{itemize}
Seeds are averaged within section first. A Mann--Whitney test was not used,
because the comparisons are paired.

The 38 comparisons form five families by question: ranking (6), against
legacy msiPL (4), encoder variants (17), segmentation (9) and S3PL (2). Holm
correction~\citep{holm1979simple} is applied within each family. With $n=8$
the smallest possible $p$ is 0.0078, so only families of at most six tests can
reach Holm significance. Effect sizes therefore carry the main weight.

A result is ``statistically supported'' only if its Holm-adjusted $p<0.05$
\emph{and} its practical label agrees. GBM is also summarised per PDX line
($n=4$, smallest $p=0.125$). Mixed models with random section and seed effects
are descriptive.

\subsection{Computation}
\label{sec:compute}
Training ran on the University's Slurm cluster (RTX 3060 and 3090 GPUs), and
ranking ran on CPU nodes.
\begin{itemize}
  \item \textbf{Caching.} The original loader re-read every batch from HDF5 on
        shared storage. Data preparation took 24.6\,s per batch against about
        0.02\,s of computation, so a 100-epoch GBM run took 11.3\,h. Reading
        each section into memory once cut it to 23\,min ($30\times$), with
        identical losses and weights. The likely cause is repeated small network
        reads, but this was not separately measured.
  \item \textbf{Reuse of saved runs.} Same-seed reruns were checked to be
        bit-identical before any saved run was reused.
  \item \textbf{Timing.} Exact ranking equality depends on the CPU model, so
        timing runs used the processor type that produced each saved ranking.
\end{itemize}
